\documentclass[a4paper,12pt]{article}
\usepackage[T2A]{fontenc}
\usepackage[utf8]{inputenc}
\usepackage[russian]{babel}
\usepackage{amsmath,amssymb,icomma}
\usepackage[a4paper,margin=22mm,headheight=15pt]{geometry}
\usepackage{microtype,booktabs,tabularx,enumitem,fancyhdr,setspace,xcolor,tikz}
\usetikzlibrary{arrows.meta,positioning,shapes.geometric,calc}
\definecolor{accent}{HTML}{008080}
\definecolor{darkaccent}{HTML}{004D4D}
\definecolor{force}{HTML}{A33B40}
\definecolor{motion}{HTML}{315D8E}
\setlength{\parindent}{0pt}
\setlength{\parskip}{5pt}
\setlist[itemize]{leftmargin=*,itemsep=1pt,topsep=2pt}
\setlist[enumerate]{leftmargin=*,itemsep=12pt,topsep=4pt}
\onehalfspacing
\pagestyle{fancy}
\fancyhf{}\fancyhead[L]{\small Динамика: баланс сил}\fancyhead[R]{\small София Лузянина}\fancyfoot[C]{\thepage}
\renewcommand{\headrulewidth}{0.3pt}
\newcommand{\fr}{F_{\text{тр}}}
\newcommand{\fc}{F_{\text{с}}}
\newcommand{\fa}{F_{\text{А}}}
\newcommand{\ft}{F_{\text{т}}}
\newcommand{\unit}[1]{\,\text{#1}}
\newcommand{\acc}{\unit{м/с$^2$}}
\newcommand{\N}{\unit{Н}}
\newcommand{\kg}{\unit{кг}}
\newcommand{\example}[1]{\par\medskip\textbf{\textcolor{darkaccent}{#1}}\par}
\tikzset{f/.style={-{Stealth[length=2.6mm]},thick,force},vel/.style={-{Stealth[length=2.6mm]},thick,motion},axis/.style={-{Stealth[length=2mm]},thin,gray!85},body/.style={draw=darkaccent,fill=accent!8,thick,rounded corners=1pt},surface/.style={thick,gray!75},flowstep/.style={draw=accent,rounded corners=2pt,align=center,text width=10.6cm,minimum height=1.05cm,inner sep=6pt},terminal/.style={ellipse,draw=accent,align=center,minimum width=4cm,minimum height=0.85cm},arrow/.style={-{Stealth[length=2.4mm]},thick,darkaccent}}
\begin{document}
\begin{titlepage}
\centering
\vspace*{28mm}
{\Large\textcolor{darkaccent}{ФИЗИКА}\par}
\vspace{14mm}
{\Huge Чек-лист\par}
\vspace{10mm}
{\LARGE\textcolor{accent}{Динамика. Второй закон Ньютона\\[2mm] и баланс сил}\par}
\vspace{14mm}
{\large Силы в механике и метод проекций\par}
\vfill
{\large София Лузянина\par}
\vspace{8mm}
27 ноября 2026 г.\par
\vspace{15mm}
{\small Лёвин Артём Александрович\\эксклюзивно для Софии Лузяниной\par}
\vspace{18mm}
\begin{tikzpicture}
\draw[accent,line width=.9pt] (0,0) .. controls (2,0.8) and (4,-0.8) .. (6,0) .. controls (8,0.8) and (10,-0.8) .. (12,0);
\end{tikzpicture}
\end{titlepage}
\setcounter{page}{1}
\section{Второй закон Ньютона}
В инерциальной системе отсчёта ускорение тела определяется векторной суммой всех сил, действующих на выбранное тело:
\[\sum_i\vec F_i=m\vec a.\]
В проекциях на две взаимно перпендикулярные оси:
\[\sum_i F_{ix}=ma_x,\qquad \sum_i F_{iy}=ma_y.\]
Если скорость постоянна по модулю и направлению, то $\vec a=0$. При торможении ускорение направлено противоположно скорости. Знак каждой проекции определяется выбранной осью.
\section{Силы: обозначения и направления}
{\renewcommand{\arraystretch}{1.2}\begin{tabularx}{\linewidth}{@{}p{4.0cm}X@{}}
\toprule Сила & Направление и формула \\\midrule
Тяжести, $m\vec g$ & Вертикально вниз; модуль $mg$.\\
Реакции опоры, $\vec N$ & Перпендикулярно поверхности, от неё к телу.\\
Натяжения нити, $\vec T$ & Вдоль натянутой нити, от тела.\\
Тяги, $\vec F_{\text{т}}$ & В заданном условием направлении.\\
Трения скольжения, $\vec F_{\text{тр}}$ & Против относительного скольжения; $\fr=\mu N$.\\
Сопротивления, $\vec F_{\text{с}}$ & Против скорости тела относительно среды.\\
Архимеда, $\vec F_{\text{А}}$ & Вертикально вверх; $\fa=\rho_{\text{ср}}gV_{\text{погр}}$.\\\bottomrule
\end{tabularx}}
\textbf{Оформление.} Над буквами векторных величин ставьте стрелки. В уравнениях проекций используйте скалярные величины со знаками. Все силы измеряются в ньютонах; масса измеряется в килограммах. $1\unit{кН}=1000\N$.
\clearpage
\section{Лифт, вес и натяжение нити}
Рассматриваем тело относительно Земли; ось $Oy$ направляем вверх. Для тела на полу лифта $N-mg=ma_y$. Для груза на вертикальной нити $T-mg=ma_y$.
\begin{center}
\begin{tikzpicture}[scale=1]
\draw[surface] (-2,-.4)--(2,-.4);\draw[body] (-.45,-.4) rectangle (.45,.4);
\draw[f] (0,0)--(0,1.55) node[above]{$\vec N$};\draw[f] (0,0)--(0,-1.45) node[below]{$m\vec g$};
\draw[axis] (-2.3,-1.3)--(-2.3,1.4) node[above]{$y$};\draw[vel] (1.45,0)--(1.45,1.1) node[above]{$\vec a$};
\begin{scope}[xshift=7cm]
\draw[surface] (-1.6,1.6)--(1.6,1.6);\draw (0,1.6)--(0,.35);\draw[body] (0,0) circle (.35);
\draw[f] (0,0)--(0,1.25) node[right]{$\vec T$};\draw[f] (0,0)--(0,-1.45) node[below]{$m\vec g$};\draw[vel] (1.45,0)--(1.45,1.1) node[above]{$\vec a$};
\end{scope}
\end{tikzpicture}
\end{center}
\begin{tabularx}{\linewidth}{@{}Xll@{}}\toprule
Состояние движения & $a_y$ & $N$ или $T$ \\\midrule
Покой или постоянная скорость & $0$ & $mg$\\
Ускорение вверх & $+a$ & $m(g+a)$\\
Ускорение вниз & $-a$ & $m(g-a)$\\
Свободное падение & $-g$ & $0$\\\bottomrule
\end{tabularx}
Вес $P$ --- сила, с которой тело действует на опору или подвес вследствие притяжения к Земле. В рассмотренных случаях $P=N$ или $P=T$ по модулю. Состояние невесомости соответствует $P=0$: сила тяжести при этом сохраняется. Формулы для контакта и натянутой нити применимы при $N\ge0$, $T\ge0$ соответственно.
\example{Пример. Лифт опускается и тормозит}
Масса пассажира $60\kg$, ускорение направлено вверх, $a=2\acc$, $g=10\acc$. Тогда
\[N=60(10+2)=720\N.\]
Показание весов увеличивается, хотя скорость направлена вниз.
\clearpage
\section{Падение и движение воздушного шара}
\begin{center}
\begin{tikzpicture}
\draw[body] (0,0) circle (.42);
\draw[f] (0,0)--(0,1.25) node[above]{$\vec F_{\text{с}}$};\draw[f] (0,0)--(0,-1.55) node[below]{$m\vec g$};
\draw[vel] (1.2,.6)--(1.2,-.6) node[below]{$\vec v$};\draw[axis] (-1.3,-1.2)--(-1.3,1.2) node[above]{$y$};
\begin{scope}[xshift=7cm]
\draw[body] (0,.15) ellipse (.65 and .85);\draw (-.27,-.63)--(-.26,-1) (.27,-.63)--(.26,-1);\draw[body] (-.3,-1.25) rectangle (.3,-1);
\draw[f] (0,.15)--(0,1.65) node[above]{$\vec F_{\text{А}}$};\draw[f] (-.2,.15)--(-.2,-1.55) node[below]{$m\vec g$};\draw[f] (.2,.15)--(.2,-1.7) node[below right]{$\vec F_{\text{с}}$};
\draw[vel] (1.5,-.4)--(1.5,.9) node[above]{$\vec v$};
\end{scope}
\end{tikzpicture}
\end{center}
\textbf{Падение тела.} При движении вниз в неподвижном воздухе сила сопротивления направлена вверх. Если выталкиванием можно пренебречь, то при оси вверх
\[\fc-mg=ma_y.\]
При равномерном падении $\fc=mg$. Если $\fc<mg$, тело ускоряется вниз; если $\fc>mg$, падающее тело тормозит.

\textbf{Подъём шара.} Для всей системы (оболочки, газа, корзины и груза) массой $m$ при движении вверх:
\[\fa-mg-\fc=ma_y,\qquad \fa=\rho_{\text{ср}}gV_{\text{погр}}.\]
Плотность $\rho_{\text{ср}}$ относится к окружающей среде. $V_{\text{погр}}$ --- объём вытесненной среды; для полностью окружённого воздухом шара это внешний вытесняющий объём системы. При спуске сопротивление направлено вверх и его проекция меняет знак.
\example{Пример. Подъём с ускорением}
Пусть $m=400\kg$, $\fa=4800\N$, $\fc=400\N$, $g=10\acc$. Тогда
\[a_y=\frac{4800-4000-400}{400}=1\acc.\]
\clearpage
\section{Горизонтальная поверхность}
Ось $Ox$ направим вправо, $Oy$ --- вверх. Брусок скользит вправо, поэтому трение направлено влево. Вертикального ускорения при сохранении контакта нет.
\begin{center}
\begin{tikzpicture}
\draw[surface] (-2,-.4)--(2,-.4);\draw[body] (-.5,-.4) rectangle (.5,.4);
\draw[f] (0,0)--(0,1.45) node[above]{$\vec N$};\draw[f] (0,0)--(0,-1.35) node[below]{$m\vec g$};\draw[f] (0,0)--(1.65,0) node[right]{$\vec F_{\text{т}}$};\draw[f] (0,0)--(-1.45,0) node[left]{$\vec F_{\text{тр}}$};
\draw[vel] (.65,1.05)--(1.65,1.05) node[right]{$\vec v$};
\begin{scope}[xshift=7.1cm]
\draw[surface] (-2,-.4)--(2,-.4);\draw[body] (-.5,-.4) rectangle (.5,.4);
\draw[f] (0,0)--(0,1.2) node[above]{$\vec N$};\draw[f] (0,0)--(0,-1.35) node[below]{$m\vec g$};\draw[f] (0,0)--(1.65,1.05) node[right]{$\vec F_{\text{т}}$};\draw[f] (0,0)--(-1.45,0) node[left]{$\vec F_{\text{тр}}$};
\draw[dashed,gray] (.5,0)--(1.7,0);\draw[gray] (.8,0) arc(0:32.5:.8);\node at (1.05,.25){$\alpha$};
\end{scope}
\end{tikzpicture}
\end{center}
\textbf{Тяга параллельна поверхности:}
\[\begin{cases}\ft-\fr=ma_x,\\N-mg=0,\\\fr=\mu N.\end{cases}\qquad a_x=\frac{\ft}{m}-\mu g.\]
\textbf{Тяга под углом $\alpha$ вверх к горизонту:}
\[\begin{cases}\ft\cos\alpha-\fr=ma_x,\\N+\ft\sin\alpha-mg=0,\\\fr=\mu N.\end{cases}\]
Отсюда $N=mg-\ft\sin\alpha$, поэтому
\[a_x=\frac{\ft(\cos\alpha+\mu\sin\alpha)}m-\mu g.\]
При толкании под углом вниз $N=mg+\ft\sin\alpha$. Трение увеличивается. Если расчёт даёт $N<0$, модель контакта с поверхностью неприменима. Закон Амонтона--Кулона в форме $\fr=\mu N$ используется для скольжения; трение покоя определяется условиями равновесия и имеет предельное значение.
\clearpage
\section{Скольжение вдоль вертикальной стены}
Пусть стена расположена слева, брусок движется вверх, сила $\ft$ направлена вверх и влево под углом $\alpha$ к вертикали. Ось $Ox$ направлена вправо, $Oy$ --- вверх.
\begin{center}
\begin{tikzpicture}
\draw[surface] (-.5,-1.9)--(-.5,2);\foreach \y in {-1.8,-1.45,...,1.7}{\draw[gray!65] (-.72,\y-.16)--(-.5,\y);}
\draw[body] (-.5,-.5) rectangle (.45,.5);
\draw[f] (0,0)--(-1.4,1.6) node[above left]{$\vec F_{\text{т}}$};
\draw[f] (0,0)--(1.7,0) node[right]{$\vec N$};\draw[f] (-.3,-.25)--(-.3,-1.5) node[left]{$\vec F_{\text{тр}}$};\draw[f] (.18,0)--(.18,-1.65) node[below]{$m\vec g$};
\draw[dashed,gray] (0,.5)--(0,1.8);\draw[gray] (0,.85) arc (90:131.2:.85);\node at (-.35,1.08){$\alpha$};
\draw[vel] (2,0.5)--(2,1.6) node[above]{$\vec v$};\draw[axis] (3,-1)--(4,-1) node[right]{$x$};\draw[axis] (3,-1)--(3,0) node[above]{$y$};
\end{tikzpicture}
\end{center}
\[\begin{cases}N-\ft\sin\alpha=0,\\\ft\cos\alpha-mg-\fr=ma_y,\\\fr=\mu N.\end{cases}\]
Следовательно,
\[N=\ft\sin\alpha,\qquad a_y=\frac{\ft(\cos\alpha-\mu\sin\alpha)}{m}-g.\]
При скольжении вниз трение направлено вверх, и знак его проекции меняется. Угол, отсчитанный от горизонтали, приводит к обмену синуса и косинуса в проекциях тяги.
\example{Пример. Подъём бруска вдоль стены}
$m=3\kg$, $\ft=60\N$, $\sin\alpha=0,6$, $\cos\alpha=0,8$, $\mu=0,25$, $g=10\acc$. Тогда $N=36\N$, $\fr=9\N$, откуда
\[a_y=\frac{48-30-9}{3}=3\acc.\]
\clearpage
\section{Наклонная плоскость}
Для наклона $\beta$ к горизонту удобно направить $Ox$ вдоль плоскости, $Oy$ --- перпендикулярно ей наружу. При высоте $h$, горизонтальном основании $b$ и длине склона $l$ имеем $\sin\beta=h/l$, $\cos\beta=b/l$, $\tg\beta=h/b$.
\begin{center}
\begin{tikzpicture}
\draw[surface] (0,0)--(6,0)--(6,3)--cycle;
\draw[gray] (.9,0) arc (0:26.565:.9);\node at (1.25,.25){$\beta$};
\begin{scope}[shift={(3,1.5)},rotate=26.565]
\draw[body] (-.45,0) rectangle (.45,.6);
\draw[f] (0,.3)--(0,1.95) node[above]{$\vec N$};\draw[f] (0,.3)--(1.45,.3) node[above right]{$\vec F_{\text{тр}}$};
\draw[f] (0,.3)--(-.8,-1.3) node[below]{$m\vec g$};
\draw[dashed,gray] (0,.3)--(-.8,.3)--(-.8,-1.3)--(0,-1.3)--cycle;
\draw[vel] (-.7,.8)--(-1.8,.8) node[above]{$\vec v$};\draw[axis] (-2.4,.15)--(-3.3,.15) node[left]{$x$};\draw[axis] (-2.4,.15)--(-2.4,1.2) node[above]{$y$};
\end{scope}
\end{tikzpicture}
\end{center}
\textbf{Спуск без тяги} (ось $Ox$ вниз по склону):
\[\begin{cases}mg\sin\beta-\fr=ma_x,\\N-mg\cos\beta=0,\\\fr=\mu N.\end{cases}\quad a_x=g(\sin\beta-\mu\cos\beta).\]
\textbf{Подъём с тягой вдоль склона} (ось $Ox$ вверх):
\[\begin{cases}\ft-mg\sin\beta-\fr=ma_x,\\N-mg\cos\beta=0,\\\fr=\mu N.\end{cases}\]
Сила тяжести всегда вертикальна. При выбранных осях модули её составляющих равны $mg\sin\beta$ вдоль и $mg\cos\beta$ поперёк плоскости. При тяге вдоль плоскости $N=mg\cos\beta$.
\example{Если тяга наклонена к поверхности}
При движении вверх и тяге под углом $\alpha$ над склоном получаем
\[N=mg\cos\beta-\ft\sin\alpha,\quad
ma_x=\ft\cos\alpha-mg\sin\beta-\mu N.\]
Этот случай объединяет два изученных приёма разложения сил. Проверяйте $N\ge0$.
\clearpage
\section{Типичные ошибки и самопроверка}
\begin{itemize}
\item Направление скорости и ускорения может различаться. При торможении ускорение имеет противоположный скорости знак.
\item Реакция опоры зависит от всех сил, действующих поперёк поверхности. Равенство $N=mg$ справедливо лишь для частного случая.
\item Трение скольжения определяется направлением относительного движения; знак его проекции выбирается по оси.
\item Сила тяжести и вес имеют разные точки приложения: первая действует на тело, второй --- на опору или подвес.
\item Угол необходимо отсчитывать от указанной оси. Проекция на эту ось равна произведению модуля силы на косинус угла с ней.
\item Уравнения по двум осям и закон трения записывайте совместно. Затем последовательно выражайте неизвестные и подставляйте числа.
\end{itemize}
\example{Пример последовательного оформления}
Брусок массой $10\kg$ тянут по горизонтальной шероховатой поверхности силой $50\N$ под углом $\alpha$ вверх, где $\sin\alpha=0,6$, $\cos\alpha=0,8$, $\mu=0,2$. Скорость направлена вправо, $g=10\acc$.
\[\begin{cases}50\cdot0,8-\fr=10a_x,\\N+50\cdot0,6-100=0,\\\fr=0,2N.\end{cases}\]
Из второго уравнения $N=70\N>0$: контакт сохраняется.
\[\fr=0,2\cdot70=14\N,\qquad 40-14=10a_x,\]
\[a_x=2,6\acc.\]
\textbf{Проверка.} Сила в числителе измеряется в ньютонах, деление на массу даёт $\text{м/с}^2$. Знак ускорения положителен; движение вправо ускоряется.
\clearpage
\section*{Алгоритм: решение задачи по динамике}
\begin{center}
\begin{tikzpicture}[node distance=8mm]
\node[terminal] (s) {Начало};
\node[flowstep,below=of s] (a) {Выберите тело, систему отсчёта и запишите данные в СИ.};
\node[flowstep,below=of a] (b) {Нарисуйте все действующие силы. Отметьте направления $\vec v$ и $\vec a$, если они известны.};
\node[flowstep,below=of b] (c) {Выберите оси вдоль движения или поверхности и перпендикулярно ей.};
\node[flowstep,below=of c] (d) {Запишите второй закон Ньютона в проекциях: $\sum F_x=ma_x$, $\sum F_y=ma_y$.};
\node[flowstep,below=of d] (e) {Добавьте законы для имеющихся сил: $\fr=\mu N$ при скольжении, $\fa=\rho gV$ при выталкивании.};
\node[flowstep,below=of e] (f) {Решите систему; проверьте знаки, единицы и условия контакта или натяжения.};
\node[terminal,below=of f] (g) {Запишите ответ};
\foreach \i/\j in {s/a,a/b,b/c,c/d,d/e,e/f,f/g}{\draw[arrow] (\i)--(\j);}
\end{tikzpicture}
\end{center}
\clearpage
\section*{Тренировочный блок}
В каждой задаче изображайте силы, выбирайте оси и записывайте уравнения в проекциях. Если требуется ускорение, указывайте его направление. Во всех задачах принимайте $g=10\acc$. Поверхности неподвижны, воздух покоится, если условием не сказано иначе.
\begin{enumerate}
\item \textbf{Движение лифта.} Пассажир стоит на напольных весах в лифте, который начинает опускаться с постоянным ускорением. Весы показывают $480\N$. На заключительном участке лифт продолжает двигаться вниз, но тормозит с ускорением такого же модуля; показание весов составляет $720\N$. Определите массу пассажира, модуль ускорения лифта и показание весов на промежуточном участке равномерного спуска. Объясните, в какую сторону направлено ускорение при торможении.
\item \textbf{Воздушный шар и балласт.} Воздушный шар вместе с оболочкой, газом, корзиной, полезным грузом и балластом имеет массу $550\kg$. Его эффективный вытесняющий объём равен $500\unit{м$^3$}$, плотность окружающего воздуха $1,2\unit{кг/м$^3$}$. Шар поднимается с постоянной скоростью. Определите архимедову силу и сопротивление воздуха. Какую массу балласта нужно сбросить, чтобы сразу после отделения шар получил ускорение $1\acc$ вверх? Считайте, что в этот момент скорость шара, объём и сопротивление воздуха не изменились, балласт отделён без толчка.
\item \textbf{Тяга и толкание.} Санки массой $20\kg$ уже скользят вправо по горизонтальной площадке. Сначала их тянут силой $100\N$ под углом $\alpha$ вверх к горизонту. Затем в другом опыте санки толкают силой того же модуля под тем же углом вниз к горизонту, сохраняя направление горизонтальной составляющей вправо. Известно, что $\sin\alpha=0,6$, $\cos\alpha=0,8$, коэффициент трения $0,25$. Для каждого случая найдите реакцию опоры, силу трения и ускорение. Во сколько раз различаются модули ускорений?
\end{enumerate}
\clearpage
\section*{Тренировочный блок (продолжение)}
\begin{enumerate}[resume]
\item \textbf{Брусок у стены.} Брусок массой $4\kg$ скользит вверх вдоль вертикальной стены, расположенной слева от него. К бруску приложена сила вверх и влево под углом $\alpha$ к вертикали ($\sin\alpha=0,6$, $\cos\alpha=0,8$). Коэффициент трения скольжения равен $0,25$. Сначала модуль силы равен $80\N$, затем его уменьшают до $40\N$. Для каждого этапа определите реакцию стены, трение и ускорение. Во втором случае рассмотрите момент, когда брусок всё ещё движется вверх. Как изменяется его скорость?
\item \textbf{Подъём груза канатом.} Груз массой $20\kg$ скользит вверх по шероховатому склону, для которого $\sin\beta=0,6$, $\cos\beta=0,8$. Канат образует угол $\alpha$ с поверхностью склона в сторону от неё ($\sin\alpha=0,6$, $\cos\alpha=0,8$). Коэффициент трения $0,25$. Найдите натяжение каната, реакцию опоры и силу трения, если ускорение направлено вверх вдоль склона и составляет $1,5\acc$. При каком натяжении реакция опоры обратится в нуль, если направление каната сохранится?
\end{enumerate}
\clearpage
\section*{Ответы}
\renewcommand{\arraystretch}{1.35}
\begin{tabularx}{\linewidth}{@{}p{.8cm}X@{}}
\toprule № & Ответ \\ \midrule
1 & $60\kg$; $2\acc$; $600\N$; при торможении ускорение вверх.\\
2 & $6000\N$; $500\N$; сбросить $50\kg$.\\
3 & Тяга вверх: $N=140\N$, $\fr=35\N$, $a=2,25\acc$ вправо. Толкание вниз: $N=260\N$, $\fr=65\N$, $a=0,75\acc$ вправо. В $3$ раза.\\
4 & При $80\N$: $N=48\N$, $\fr=12\N$, $a=3\acc$ вверх. При $40\N$: $N=24\N$, $\fr=6\N$, $a=3,5\acc$ вниз; скорость по модулю уменьшается.\\
5 & $T=200\N$; $N=40\N$; $\fr=10\N$; граница контакта $T=800/3\N\approx266,7\N$.\\
\bottomrule
\end{tabularx}
\vspace{7mm}
\textcolor{darkaccent}{\textbf{Итоговая памятка.}} Сначала определите силы и их направления, затем составьте уравнения проекций. Реакцию опоры вычисляйте из уравнения по нормали к поверхности. Сохраняйте знаки ускорения до окончания расчёта.
\end{document}
